\section{Slutsatser och avslutande diskussion}

När jag ser tillbaka på min LIA-period på Insighta känner jag att jag har utvecklats mycket, både tekniskt och personligt. I inledningen satte jag upp mål om att få praktisk erfarenhet av DevOps-arbete, lära mig Kubernetes, Terraform och GCP samt förstå hur ett riktigt team arbetar tillsammans. Jag tycker att jag har nått dessa mål i hög grad.

Det viktigaste jag tar med mig är att jag har fått självförtroende i att jag faktiskt kan arbeta i en professionell miljö. Jag känner att jag kan ta ansvar, lösa uppgifter och bli betrodd med riktiga projekt. Detta är något som betyder mycket för mig, eftersom jag inför LIA:n var osäker på hur väl jag skulle klara av avancerade verktyg och arbetsflöden. Nu vet jag att jag kan.

Jag har också fått en djupare förståelse för DevOps-arbetssättet. Jag lärde mig att arbeta agilt, följa sprintar, använda tickets och göra pull requests. Detta gav mig en tydligare bild av hur team organiserar sitt arbete och hur viktigt det är med kommunikation och tydlighet. Jag har dessutom fått erfarenhet av flera verktyg som vi bara pratat om i skolan, men inte använt praktiskt, som Kubernetes, GKE, Terraform, CDKTF och avancerad logghantering genom GCP. Genom att arbeta med dem i riktiga projekt förstår jag nu mycket bättre hur de fungerar.

Det som jag lärde mig allra mest av var att deploya Dagster och Postgres till Kubernetes, sätta upp GKE-kluster och arbeta med Terraform. Jag fick arbeta med komplexa problem som krävde felsökning och tålamod. Genom detta lärde jag mig att inte stressa upp mig, utan testa en sak i taget tills jag hittade lösningen.

Jag känner också att min förståelse för molninfrastruktur har blivit mycket bredare. Nu när jag har jobbat praktiskt med GCP och Kubernetes förstår jag hur viktigt det är för framtida DevOps- och MLOps-roller.

Slutligen har denna LIA gjort mig ännu mer motiverad. Jag planerar att fortsätta på samma företag i min andra LIA, och efter den siktar jag på att ta certifieringar, både CNCF Kubernetes-certifiering och en GCP-certifiering. Mitt mål är att göra dessa vid slutet av min andra LIA-period.

Sammanfattningsvis har min LIA på Insighta varit en mycket värdefull erfarenhet som gjort mig tryggare i min yrkesroll och mer säker på att jag vill fortsätta arbeta inom DevOps och molnteknologi.